\section{总图：十个分支与一张依赖网}

\subsection{状态图例}

全篇每个条目后面都跟一个状态词，取值为：

\begin{center}
\begin{tabular}{L{2.6cm}L{12.4cm}}
\toprule
状态 & 含义 \\
\midrule
已实现 & 源码头文件标注「已实现」，函数体有真实算法 \\
框架 & 签名与注释已定、函数体全部 \texttt{throw}；注释即规格，公式可信，数值未验证 \\
骨架 & 同上，但源码自述为「骨架」；公式仍是规格 \\
未建 & 只在 \texttt{NOW.md} 里规划，文件不存在 \\
缺位 & 库里既没有位置、也没有规划；必须在第 \ref{sec:audit} 节的缺口里登记 \\
\bottomrule
\end{tabular}
\end{center}

\textbf{框架与骨架的数学含量并不低}：它们是「先写契约后写代码」的产物，公式段往往是全仓最完整的。
但它们的\textbf{判据全部未经验证}——条目里凡是标「未验证」的阈值，都只代表意图，不代表实测。

\subsection{十个分支与条目编号}

条目编号是全局的、不按章的先后。同一编号在全文中唯一，交叉引用都用它。

\begin{center}
\begin{tabular}{L{2.5cm}L{2.9cm}L{8.4cm}}
\toprule
分支 & 编号 & 覆盖对象 \\
\midrule
分析基础 & A1--A6 & 机器精度、累加误差、前向/后向误差、条件数、收敛阶、相消 \\
线性代数 & L1--L10 & 线性映射、矩阵乘、转置与伴随、内积、范数与迹、谱、奇异值、正交相似 \\
数值线性代数 & N1--N13 & 部分主元、Householder、Givens、MGS、Cholesky、Padé、Jacobi、$A^{H}A$ 路线、秩判定 \\
复分析 & C1--C6 & 复代数、Smith 除法、欧拉公式、主值与分支切割、开方、幂 \\
调和分析 & H1--H11 & DFT、归一化、卷积定理、频轴、radix-2、Parseval、级数投影、Gibbs、谱收敛 \\
概率 & P1--P10 & 特征函数、矩、接缝、经验特征函数、Box--Muller、erf/$\Phi$、分位数 \\
微分与自动微分 & D1--D8 & 对偶数、链式法则、雅可比布局、种子 sweep、反向模式、Lipschitz、RK 阶、辛结构 \\
几何与李群 & G1--G10 & 叉积与 hat、Rodrigues、李括号、四元数代数、双覆盖、插值、Shepperd \\
离散化与变分 & F1--F9 & 网格与自由度、P1 刚度、弱形式、一致载荷、Dirichlet、收敛阶、CSR、CG \\
张量代数 & T1--T7 & 形状与步长、偏移、连续性、零拷贝视图、reshape 合法性、广播 \\
微分形式 & M1--M7 & $k$-形式、外积、外微分、拉回、Hodge、Stokes、$n=3$ 字典 \\
\bottomrule
\end{tabular}
\end{center}

\subsection{依赖网（ASCII）}

箭头读作「是……的基础」。方括号里是状态。

\begin{verbatim}
 A1 eps --> A2 8n*eps --> A3 fwd/bwd err --> A4 cond number
   |            |               |                |
   |            v               v                v
   |          N1 pivot     N5 MGS ill-cond   N12 A^H A squared
   v            |                                ^
 C2 Smith       v                                |
            N3 Householder --> N6 thin QR    L7 singular values
                 |                 |              ^
                 v                 v              |
               N4 Givens      N6b least sq   N11 Jacobi spectrum

 L1 row-major --> L2 matmul --> L3 transpose/adjoint
   |                              |
   v                              v
 T1 shape/stride --> T2 offset --> T3 contiguous --> T4 view
                                                        |
                                                        v
                                                   T5 reshape
 D1 dual eps^2=0 --> D2 chain rule --> D3 jacobian layout
                                        |--> D4 seed sweep (col)
                                        \--> D5 reverse mode (row)

 H1 DFT --> H2 1/N --> H3 circular conv --> H4 linear conv (pad)
   |          |                                |
   v          v                                v
 H5 freq axis  H8 Parseval            H10 Gibbs / H11 spectral decay
   |
   v
 P1 char func (+i, opposite sign) --> P3 conv theorem --> P6 empirical
                                          |
                                          v
                                     P4 moments = derivatives

 G1 cross/hat --> G2 Rodrigues --> G4 Lie bracket
      |               |
      v               v
 G6 quat double cover --> G8 slerp --> G9 Shepperd

 D6 Lipschitz --> D7 Euler/RK4 order --> D8 symplectic/Verlet

 F1 grid --> F2 P1 stiffness --> F3 weak form --> F4 load
                                   |               |
                                   v               v
                               F5 Dirichlet --> F6 order 2
                                   |
                                   v
                              F7 CSR --> F8 CG --> F9 convergence

 M3 d^2=0 --> M6 Stokes --> M7 grad/curl/div  (same root as F3)
\end{verbatim}

\subsection{读法：三层}

\begin{enumerate}[label=（\arabic*）]
  \item \textbf{查漏}：翻到某分支，逐条看「常见误解」——能反驳自己直觉的，才说明真懂；
  \item \textbf{补基础}：条目末的「前置」给出向上追溯的路径，追到 A 开头的条目为止；
  \item \textbf{建直觉}：星号节「串讲」是跨分支的边，专门讲「为什么这两个看起来无关的东西是同一件事」。
\end{enumerate}

本份清单共 \textbf{97 条}知识点：已实现 \textbf{49 条}、框架 \textbf{28 条}、骨架 \textbf{3 条}、
未建 \textbf{16 条}、缺位 \textbf{1 条}。这五类数字的差别比总数重要得多，含义见第 \ref{sec:audit} 节。

状态词是\textbf{白名单五值}、可脚本核对（\texttt{math/survey/verify.py} 的 V2）：
其中「未建」指只在 \texttt{NOW.md} 里规划、「缺位」指库里既无位置也无规划。
两个词不许混用，混了计数就对不上。
